\documentclass[12pt]{article}

\usepackage[utf8]{inputenc}
\usepackage[T1]{fontenc}
\usepackage[english]{babel}

\usepackage{amsmath, amssymb}
\usepackage{physics}
\usepackage{geometry}
\usepackage{graphicx}
\usepackage{setspace}
\usepackage{hyperref}

\geometry{margin=2.5cm}
\onehalfspacing

\title{Longitudinal Waves in Coupled Thermoelasticity}
\author{}
\date{}

\begin{document}

\maketitle

\section{Introduction}

\subsection{General Background}

Geological media comprise natural solid Earth materials such as sedimentary, igneous, and metamorphic rocks, together with fractured, weathered, porous, and stress-bearing rock masses derived from them. From a geological point of view, these materials are products of deposition, diagenesis, magmatism, metamorphism, deformation, erosion, and surface alteration. From a physical point of view, they are solid media with complex internal architecture, whose behavior is controlled not by a single property but by the combined influence of mineral composition, density, porosity, grain arrangement, fracture systems, elastic stiffness, temperature, and in situ stress conditions.

Rock physics studies the relationship between geophysical observations and the physical properties of rocks, while geological analysis explains why these properties vary between lithologies, depositional settings, burial histories, and structural environments. Therefore, a geological medium cannot be adequately represented as an ideal uniform solid without acknowledging that its present response reflects both its physical state and its geological history (Mavko et al., 2009; Jaeger et al., 2007; Turcotte and Schubert, 2014).

This dependence on internal structure is especially important because rocks belonging to the same broad lithological class may behave differently under different geological and thermodynamic conditions. Thermal properties of rocks and minerals vary with composition, porosity, water content, pressure, and temperature. Similarly, seismic velocities and elastic moduli depend on porosity, differential pressure, degree of consolidation, lithology, fracture spacing, weathering state, and the arrangement of mineral grains. In competent bedrock, fracture spacing and crack orientation may influence shear-wave velocity more strongly than hardness alone. In unconsolidated or weakly consolidated sediments, grain-size distribution, packing, cementation, and effective stress strongly affect stiffness and wave velocity.

Mineral composition is also significant because minerals differ in density, elastic constants, thermal conductivity, heat capacity, and thermal expansion. For this reason, the behavior of geological media must be understood as the combined result of composition, microstructure, and in situ conditions rather than as an intrinsic constant of ``rock'' in the abstract (Mavko et al., 2009; Robertson, 1988; Jaeger et al., 2007).

\subsection{Physical and Geological Motivation}

Within geophysics, the propagation of waves in rocks is one of the principal ways by which the subsurface becomes observable. Seismic and acoustic methods rely on elastic disturbances traveling through the Earth, and their interpretation depends on the relationship between wave behavior and rock properties. Travel times, amplitudes, reflected arrivals, refracted waves, and changes in waveform are used to infer layer thickness, depth to bedrock, lithological contacts, stratigraphic boundaries, fracture zones, and lateral heterogeneity. In this sense, wave propagation forms the physical basis of seismic interpretation and a large part of modern subsurface characterization.

The same principles are relevant in engineering geology, where geophysical surveys are used to support site investigation, reduce uncertainty between boreholes, and provide continuity where direct drilling or excavation is limited. Wave propagation in geological media is therefore not only a subject of theoretical mechanics, but also a foundation for practical geological reasoning about the structure and condition of the subsurface (Aki and Richards, 2002; Mavko et al., 2009).

The importance of this subject becomes even clearer in applications related to geological hazards and energy resources. Near-surface shear-wave velocity is known to influence the amplification and duration of earthquake ground motion, meaning that the elastic properties of rock and sediment directly affect seismic hazard evaluation and engineering response. At larger scales, wave propagation is central to understanding crustal structure, earthquake source processes, and the mechanical state of geological formations.

Geothermal systems provide another important motivation. Geothermal energy is stored mainly in hot rocks, but its extraction depends on the transfer of heat through the rock mass and, in many cases, the movement of fluids through fractures and pores. In such environments, temperature is not merely a background variable. It can affect deformation, cracking, permeability, stiffness, and the propagation of mechanical disturbances. Therefore, the study of waves in thermally affected geological media is relevant to geophysics, engineering geology, geothermal resource assessment, and geological hazard analysis (Turcotte and Schubert, 2014; Jaeger et al., 2007).

\subsection{Thermal--Mechanical Interaction in Rocks}

A central physical issue in this context is the interaction between temperature and deformation. In classical thermoelasticity, an increase in temperature produces thermal expansion, whereas a decrease in temperature produces thermal contraction. If a body were completely free to deform, the resulting strain would be governed mainly by the thermal expansion coefficient and the temperature change. Geological media, however, are rarely free bodies. Rock layers are confined by surrounding strata, mineral grains are constrained by neighboring grains, and rock masses are bounded by faults, joints, lithological contrasts, and regional stress fields.

Under such conditions, thermal expansion or contraction may generate thermal stresses. Heating under constraint commonly produces compressive stress, whereas cooling may promote tensile stress. At the mineral scale, different crystals may expand by different amounts and sometimes in different directions, so temperature changes may generate local stress concentrations and microcracks even when the average temperature change is moderate. These processes are geologically important because repeated heating and cooling can modify stiffness, crack density, permeability, and the long-term integrity of rock masses (Nowacki, 1986; Boley and Weiner, 1960; Robertson, 1988).

The geological significance of thermal stress is reinforced by the fact that temperature changes are capable of producing visible structural effects in rocks. Cooling contraction is a recognized cause of jointing in igneous rocks, and unequal thermal expansion between mineral grains may contribute to internal cracking within polymineralic aggregates. More generally, thermal cycling can produce progressive rock damage and alter both mechanical and transport properties. This means that temperature is not simply an external field imposed on an otherwise mechanical problem. It can change the effective stiffness of a rock mass, influence fracture development, modify pathways for heat transfer, and thereby affect subsequent wave propagation.

In geological environments where thermal gradients are strong or time-dependent, such as volcanic regions, geothermal reservoirs, shallow weathering zones, or rocks subjected to repeated heating and cooling, thermoelastic effects become part of the natural physics of the medium itself (Turcotte and Schubert, 2014; Jaeger et al., 2007; Robertson, 1988).

\subsection{Heat Conduction and Elastic Wave Propagation}

A further distinction must be made between thermal and mechanical disturbances. In the classical theory considered here, temperature evolves primarily through heat conduction, which is governed by thermal conductivity and heat capacity. These properties are commonly combined in thermal diffusivity, which controls how rapidly a material adjusts to a change in temperature. Materials with high thermal diffusivity respond more quickly to thermal disturbances than materials with low thermal diffusivity.

Mechanical disturbances, in contrast, propagate as elastic waves whose velocities depend mainly on elastic moduli and density. Thermoelasticity links these two processes by recognizing that temperature changes generate strain and stress, while deformation may also enter the thermal energy balance of the solid. Conceptually, thermoelastic wave propagation in a rock therefore involves the coupled evolution of temperature, displacement, strain, and stress rather than two unrelated problems placed side by side (Nowacki, 1986; Boley and Weiner, 1960).

The elastic-wave part of this coupled problem is usually introduced through the basic body waves of solid media. Compressional waves, or P-waves, are associated mainly with local volume change, while shear waves, or S-waves, are associated mainly with change of shape. In a homogeneous elastic medium, P-waves travel faster than S-waves, and the velocities of both depend on density and elastic stiffness. When waves encounter contrasts in these properties, part of the wave energy is transmitted and part is reflected. This is why elastic waves can be used to infer subsurface structure.

\subsection{Directional Behavior in Geological Wave Propagation}

In real geological media, wave velocities are not controlled only by density and elastic moduli in an ideal uniform sense. Layering, foliation, aligned minerals, pores, joints, faults, and fracture systems may make the response dependent on direction. Consequently, the same rock mass may transmit waves differently along bedding, across bedding, parallel to fractures, or perpendicular to fractures. This directional behavior is a natural geological consequence of anisotropy, heterogeneity, and internal structure rather than a purely mathematical complication (Aki and Richards, 2002; Mavko et al., 2009).

Directional behavior is especially important for geological wave propagation because rocks rarely possess the same physical properties in all directions. Sedimentary rocks may contain bedding planes, laminations, preferred grain orientations, and compaction-related fabric. Metamorphic rocks may contain foliation, schistosity, or mineral alignment. Igneous and crystalline rocks may contain cooling joints, microcracks, and fracture networks. In each case, the direction of wave travel relative to these structures can influence velocity, attenuation, scattering, and the partition of wave energy.

A wave traveling parallel to a continuous stiff layer may propagate differently from a wave crossing alternating layers with contrasting stiffness. Similarly, a wave moving parallel to open fractures may experience a different effective stiffness than a wave moving normal to them. Temperature can further affect this directional response by changing stress conditions, opening or closing cracks, and producing thermal expansion that may be constrained differently along different structural directions. Thus, the directional behavior of thermoelastic wave propagation is closely related to the geological fabric of the medium and to the thermal--mechanical state in which the wave travels.

\subsection{Research Gap}

Despite the established importance of elastic waves, heat conduction, and thermal stress as separate topics, their combined interpretation in geological media remains a significant theoretical challenge. Classical elasticity provides a basis for describing displacement, strain, and stress, while heat-conduction theory describes the evolution of temperature. Classical thermoelasticity connects these fields through thermal expansion and thermal stress. However, geological materials introduce additional complexity because their properties are heterogeneous, direction-dependent, scale-dependent, and sensitive to temperature, pressure, porosity, and fractures.

As a result, there is a gap between idealized thermoelastic formulations and the physical behavior of real geological media. In particular, a clear theoretical explanation is needed of how heat conduction, thermal expansion, elastic deformation, thermal stress, and directional wave behavior are connected when the medium is a rock rather than an abstract homogeneous solid. This gap does not imply that classical theory is insufficient as a foundation. Rather, it shows the need to present the classical framework in a geological context, where material properties and directional structure are physically meaningful.

\subsection{Aim and Objectives of the Theoretical Part}

The aim of the theoretical part of this thesis is to establish the physical and geological basis required to understand thermoelastic wave propagation in geological media.

To achieve this aim, the following objectives are defined:

\begin{itemize}
    \item to describe the main geological properties that control the mechanical and thermal behavior of rocks, including density, porosity, mineral composition, elastic stiffness, thermal conductivity, heat capacity, and thermal expansion;

    \item to explain the classical concepts of elasticity, strain, stress, and elastic wave propagation in solid media;

    \item to introduce heat conduction and the role of thermal diffusivity in the evolution of temperature fields;

    \item to explain how thermal expansion under geological constraint produces thermal stress;

    \item to connect these concepts to the directional behavior of waves in structured geological media.
\end{itemize}

These objectives are theoretical in nature and do not require the assumption of experimental results.

\subsection{Structure of the Theoretical Discussion}

Against this background, the physical problem considered in this thesis may be stated in classical terms. Given a geological material with density, elastic properties, coefficient of thermal expansion, thermal conductivity, and heat capacity, the task is to describe how a thermal disturbance and the associated mechanical deformation evolve within the medium. More specifically, one seeks the time-dependent temperature field together with the corresponding displacement, strain, and stress fields, while accounting for the fact that temperature changes generate thermal strains and that these strains alter the mechanical state of the solid.

In classical thermoelastic theory, this description is constructed from the equations of elasticity, the constitutive law including thermal strain, and the heat-conduction equation with appropriate initial and boundary conditions. For rocks, this formulation provides the essential starting point from which more geologically realistic situations, including heterogeneity, fracturing, anisotropy, and stress dependence, can be interpreted (Nowacki, 1986; Boley and Weiner, 1960).

The purpose of the present Introduction is therefore to motivate the physical-geological problem and to define the conceptual framework for the following chapter. Chapter 1 develops the relevant foundations in a systematic manner: geological material properties, elasticity, thermal expansion, heat conduction, thermoelastic stress, and elastic wave propagation. Later parts of the thesis may build upon this theoretical basis, but the role of the Introduction and Chapter 1 is limited to the classical physical and geological concepts needed to understand how temperature and deformation interact during wave propagation in rocks. This restriction is important because it keeps the theoretical foundation focused on the natural behavior of geological media and on the governing principles of thermoelasticity.

\end{document}